% Exported from Math++. Compile with pdflatex.
\documentclass[11pt,a4paper]{article}
\usepackage[T1]{fontenc}
\usepackage[utf8]{inputenc}
\usepackage{lmodern}
\usepackage{amsmath,amssymb,amsthm,mathtools}
\usepackage{graphicx,booktabs,xcolor}
\usepackage[hidelinks]{hyperref}
\usepackage[margin=25mm]{geometry}
\providecommand{\degree}{\ensuremath{^\circ}}
\providecommand{\xmark}{\ensuremath{\times}}
\providecommand{\omicron}{o}
\providecommand{\oiint}{\mathop{\oint\mkern-12mu\oint}}
\theoremstyle{plain}
\newtheorem{theorem}{Theorem}
\newtheorem{lemma}[theorem]{Lemma}
\newtheorem{corollary}[theorem]{Corollary}
\newtheorem{proposition}[theorem]{Proposition}
\newtheorem{conjecture}[theorem]{Conjecture}
\theoremstyle{definition}
\newtheorem{definition}[theorem]{Definition}
\newtheorem{example}[theorem]{Example}
\theoremstyle{remark}
\newtheorem*{remark}{Remark}
\setlength{\parindent}{0pt}\setlength{\parskip}{0.6em}
\title{Checking identities and answers}
\date{}
\begin{document}
\maketitle
Write \texttt{\textbackslash{}overset\{?\}\{=\}} (or \texttt{?}, or \texttt{?=}) between two sides to check them. Math++ tries a
symbolic proof first, then compares the sides at random points, and marks the line $\checkmark$ or $\xmark$;
a $\xmark$ shows a point where the sides differ.

\section*{Identities}

\begin{equation*}
  \sin^2 x + \cos^2 x \overset{\checkmark }{=} 1
\end{equation*}
\begin{equation*}
  \sin 2x \overset{\checkmark }{=} 2 \sin x \cos x
\end{equation*}
\begin{equation*}
  \cos 2x \overset{\checkmark }{=} \cos^2 x - \sin^2 x
\end{equation*}
\begin{equation*}
  \tan^2 x + 1 \overset{\checkmark }{=} \sec^2 x
\end{equation*}
\begin{equation*}
  \sin(x + y) \overset{\checkmark }{=} \sin x \cos y + \cos x \sin y
\end{equation*}
\begin{equation*}
  \cosh^2 x - \sinh^2 x \checkmark  1
\end{equation*}
\begin{equation*}
  (a + b)^2 \checkmark = a^2 + 2ab + b^2
\end{equation*}
\section*{Homework answers}

\begin{equation*}
  \frac{d}{dx} x^2 \sin x \overset{\checkmark }{=} 2x \sin x + x^2 \cos x
\end{equation*}
\begin{equation*}
  \int x e^x \, dx \overset{\checkmark }{=} (x - 1) e^x
\end{equation*}
An antiderivative may differ from Math++'s by a constant and still be right:

\begin{equation*}
  \int \sin x \cos x \, dx \overset{\checkmark }{=} -\frac{\cos^2 x}{2}
\end{equation*}
\begin{equation*}
  \lim_{x \to 0} \frac{1 - \cos x}{x^2} \overset{\checkmark }{=} \frac{1}{2}
\end{equation*}
\begin{equation*}
  \sum_{n=1}^{\infty} \frac{1}{2^n} \overset{\checkmark }{=} 1
\end{equation*}
\section*{Mistakes}

\begin{equation*}
  (x + 1)^2 \overset{\xmark }{=} x^2 + 1
\end{equation*}
\begin{equation*}
  \sqrt{x^2} \overset{\xmark }{=} x
\end{equation*}
\begin{equation*}
  \frac{d}{dx} \ln(2x) \overset{\xmark }{=} \frac{2}{x}
\end{equation*}
\begin{equation*}
  \int_0^1 x^2 \, dx \overset{\xmark }{=} \frac{1}{2}
\end{equation*}
\end{document}
